\section{面向对象（下）}

\begin{question}
	在Java中，以下哪个关键字用于实现继承？
	\begin{options}
		\item extends
		\item implements
		\item interface
		\item abstract
	\end{options}
	\begin{answer}
		A
	\end{answer}
	\begin{explanation}
		\texttt{extends} 用于类继承；\texttt{implements} 用于实现接口。
	\end{explanation}
\end{question}

\begin{question}
	以下关于抽象类的说法错误的是
	\begin{options}
		\item 抽象类可以有抽象方法
		\item 抽象类不能被实例化
		\item 抽象类的子类必须实现所有抽象方法
		\item 抽象类中可以有非抽象方法
	\end{options}
	\begin{answer}
		C
	\end{answer}
	\begin{explanation}
		如果子类也是抽象类，则可以不实现父类的抽象方法。
	\end{explanation}
\end{question}

\begin{question}
	以下哪个关键字用于定义接口？
	\begin{options}
		\item interface
		\item abstract class
		\item class
		\item extends
	\end{options}
	\begin{answer}
		A
	\end{answer}
	\begin{explanation}
		接口使用 \texttt{interface} 关键字定义。
	\end{explanation}
\end{question}

\begin{question}
	以下关于接口的说法错误的是
	\begin{options}
		\item 接口中的方法默认是public abstract修饰的
		\item 接口中的变量默认是public static final修饰的
		\item 一个类可以实现多个接口
		\item 接口可以继承类
	\end{options}
	\begin{answer}
		D
	\end{answer}
	\begin{explanation}
		接口只能继承接口，不能继承类。
	\end{explanation}
\end{question}

\begin{question}
	以下哪个关键字用于修饰内部类，使其可以访问外部类的私有成员？
	\begin{options}
		\item static
		\item final
		\item private
		\item public
	\end{options}
	\begin{answer}
		A
	\end{answer}
	\begin{explanation}
		静态内部类可以访问外部类的私有静态成员；非静态内部类天然可访问外部类所有成员（包括私有）。
	\end{explanation}
\end{question}

\begin{question}
	以下关于匿名内部类的说法错误的是
	\begin{options}
		\item 没有类名
		\item 可以继承其他类或实现接口
		\item 只能使用一次
		\item 不能定义构造方法
	\end{options}
	\begin{answer}
		C
	\end{answer}
	\begin{explanation}
		匿名内部类虽然没有显式类名，但可以在创建时多次使用（例如多次 new），只是每次都是新类。
	\end{explanation}
\end{question}

\begin{question}
	在Java中，以下哪个方法用于比较两个对象是否相等？
	\begin{options}
		\item equals()
		\item compareTo()
		\item isEqual()
		\item same()
	\end{options}
	\begin{answer}
		A
	\end{answer}
	\begin{explanation}
		\texttt{equals()} 用于逻辑相等比较。
	\end{explanation}
\end{question}

\begin{question}
	以下关于Java中对象克隆的说法错误的是
	\begin{options}
		\item 可以通过实现Cloneable接口来实现对象克隆
		\item clone()方法是在Object类中定义的
		\item 克隆分为浅克隆和深克隆
		\item 默认的clone()方法实现的是深克隆
	\end{options}
	\begin{answer}
		D
	\end{answer}
	\begin{explanation}
		默认 \texttt{clone()} 是浅克隆。
	\end{explanation}
\end{question}
